\documentclass[a4j, 9pt]{ltjsarticle}

\usepackage{../mypackage}

\begin{document}

\section{論理学}

  数学では様々な議題を論証的に議論する。これには、\textbf{証明}が不可欠である。
  例えば次の主張

  \begin{equation*}
    \begin{split}
      「幅$\ds a$で高さ$\ds b$の長方形$\ds A$ & の面積$\ds S_A$と\\
                                            & 同様の幅と高さを持つ三角形$\ds B$の面積$\ds S_B$の間には$\ds S_A = 2 \times S_B$なる関係が成立する」
    \end{split}
  \end{equation*}

  これは数学的に証明される主張であるが、そのためには幾何学的に\textbf{長方形}と\textbf{三角形}を定義し、さらにその幾何量に合うように解析的に\textbf{面積}が定義されていて初めて証明に至る。\par

  このような、微妙であり、重要な論理関係を明瞭に考えるには\textbf{論理学}という少々哲学的なところから議論をする必要がある。
  そこで本セクションでは数学的議論をするために必要な論理学に立ち入ってみることにしよう。
  なお、本セクション以降では論理関係を明確化するため論理記号等を必要に応じて用いることにする。

  \subsection{命題と条件}

    \subsubsection{命題}
      論理学が興味を持つのは、複雑な主張においても必ずその主張が"正しい"、あるいは"誤りである"と推論を打ち立てられるような主張(statement)である。
      このように"正しい"または"誤りである"のどちらかが定められる・定めうる主張のことを論理学では\textbf{命題}(proposition)と言う。
      また、主張が正しいことを\textbf{真}と言い、主張が誤りであることを\textbf{偽}と言う。

    \subsubsection{条件}
      次の数学計算式があるとする。
      
      \begin{equation*}
        \begin{split}
          $\ds x + 3 = 0$
        \end{split}
      \end{equation*}

      この式は、数学的には方程式というが、論理学的にこの主張を見てやると$\ds x$という文字(数学的には変数と言う)の値次第で真偽が変わってしまう主張である。
      例えば$\ds x = -3$のときこの式が成立するので真であり、この値をこの方程式の解というが、$\ds x = 1$のときこの式は成立しないので偽である。
      もちろんこれも命題の一つである事には変わりがない。
      このような文字を論理学では\textbf{変項}と呼ぶ。\par
      また、このような変項$\ds x$の値次第で真偽が変わってしまう命題のことを\textbf{変項$\ds x$についての条件}と言う。
      \footnote{
        次のセクションで説明する集合の記号を用いると$\ds 条件 \subset 命題$である。
      }

  \subsection{命題論理と述語論理}
    19世紀に入る頃、論証を支える論理そのものも計算式的に記述できることがわかってきた。
    つまるところ、論証は演算と言っても過言ではない。\par
    例えば次のような数学計算式は

    \begin{equation*}
      \begin{split}
        $\ds 2 + 3 = 5$
      \end{split}
    \end{equation*}

    数である$\ds 2$と$\ds 3$を\textbf{作用する対象}とし、これらの間に加法\textbf{演算子}$\ds +$を組み込むことで$\ds 2 + 3$という演算を行い、その結果が$\ds 5$という数と等しいという\textbf{主張}を表す。\par
    これに対応するように、論証においては\textbf{作用する対象を命題}とし、\textbf{演算子を"かつ"(and)や"または"(or)など}と対応させる。\par
    例として次のような主張を考えてみる。

    \begin{equation*}
      \begin{split}
        ワンちゃんはかわいい　かつ　ワンちゃんは賢い
      \end{split}
    \end{equation*}

    \reviewstrike{無論これは事実である。}
    この主張は、「ワンちゃんはかわいい」という命題と、「ワンちゃんは賢い」という命題が作用する対象の命題であり、演算子が「かつ」である。\par
    一般に、このような\textbf{複数の命題を演算子で合成した命題}を\textbf{合成命題}という。
    また、「ワンちゃんはかわいい」のような合成命題を構成する基本命題をその合成命題の\textbf{命題変項}または\textbf{命題変数}と呼ぶ。
    このように、命題に対して「かつ」のような演算子が作用して合成命題を成していることから論理式という言葉も納得である。

    \subsubsection{命題論理}
      今見たように、「かつ」のような演算子がある。
      このような演算子を次に示すように記号を定義する。\par

      \begin{center}
        \begin{table}[hbtp]
          \centering
          \begin{tabular}{lcr}
            \hline
            論理記号  & 英語文  &  日本語文  \\
            \hline \hline
              $\ds\land$                              & and         & かつ\\
              $\ds\lor$                               & or          & または\\
              $\ds\Longrightarrow, \longrightarrow$   & if- then... & ならば\\
              $\ds\lnot, \bar{　}$                    & not         & でない(否定)\\
          \end{tabular}
        \end{table}
      \end{center}

      これらの演算子を\textbf{論理演算子}、ないしは\textbf{論理結合子}と呼ぶ。
      これらの演算子は論理演算において最も低いレベルの演算子であり、これらを用いて構成される合成命題を\textbf{命題論理}とも呼ぶ。\par
      最も低いレベルと言っているのは、数学において例えば$\ds 2 \times 3 = 6$という乗法は$\ds 2 + 2 + 2 = 6$のようにして加法を用いて分解可能であることのような、
      より低い程度の演算で書き換え可能であることと同様に、任意の合成命題はこの命題論理に分解可能であるということである。
      \footnote{
        後に述べるド・モルガン則によれば、命題$\ds p, q$に対する$\ds p \land q$も$\ds \lnot(\lnot{p} \lor \lnot{q})$と書くことができるため、
        より低レベルな命題論理としては$\ds \land$を省くと考えられるが、ここでは簡単のため$\ds \land$も含むものとする。ちなみにこれは後に述べるNANDの否定である。
      }
      \par

      今、第一のレベルの論理演算子を定義したが、これらは言語上曖昧な表現が多いために誤解を防ぐために定義したのである。
      ゆえにこれらを用いて記述される論理式の意味が曖昧なままでは演算子を定義した意味がない。\par
      そこで次の表に与えるように、論理演算子の意味を定義する。

      \begin{center}
        \begin{table}[h]
          \centering
          \begin{tabular}{c|c||c|c|c|c}
            \truthtable{A,B}{$A$,$B$}{A&B, A|B, >>(A;B), !A} {$A \land B, A \lor B, A \Rightarrow B, \lnot A$}{$T$}{$F$}
          \end{tabular}
        \end{table}
      \end{center}
      
      Tが真(true)、Fが偽(false)である。
      この表のことを\textbf{真理表}と言う。\par

      \paragraph{条件記号の曖昧さ}
        数学では"かつ"を表したり、場合分けをするときにしばしば"\{"を用いる。
        例えば$\ds x + y = 3$と$\ds x - y = 5$の2元1次連立方程式は

        \begin{center}
          \begin{cases}
            $\ds x + y = 3$\\
            $\ds x - y = 5$\\
          \end{cases}
        \end{center}

        と書くが、これが表す論理式は

        \begin{equation*}
          \begin{split}
            $\ds x + y = 3 \land x - y = 5$
          \end{split}
        \end{equation*}

        である。\par
        一方で$\ds x = |a|$の解は$\ds a$の符号によって変わるため

        \begin{center}
          $\ds x =$
          \begin{cases}
            $\ds a$   & if　$\ds a \geq 0$\\
            $\ds -a$  & if　$\ds a < 0$\\
          \end{cases}
        \end{center}
        \footnote{
          あるいは$\ds \pm$や$\ds \mp$を用いて$\ds x = \pm a$とも書く。
        }

        のようにして書き、これが表す論理式は

        \begin{equation*}
          \begin{split}
            $\ds x = a \lor x = -a$
          \end{split}
        \end{equation*}
        である。
        このように"\{"の表す論理式は、文脈によって変わってしまい少々曖昧であるので、本来は論理式で書くのがより確実である。

      \columnbreak

      $\ds p, q$を命題変項として、先ほど定義した論理演算子を用いて次のような演算子が定義できる。

      \begin{center}
        \begin{table}[hbtp]
          \centering
          \begin{tabular}{lcr}
            \hline
            論理記号  & 英語文での意味  &  定義内容  \\
            \hline \hline
              $\ds p \Longleftrightarrow q$ & $\ds p$ iff (if and only if) $\ds q$ & $\ds "p \Rightarrow q" \land "p \Leftarrow q"$\\
          \end{tabular}
        \end{table}
      \end{center}
      
      これは$\ds p$と$\ds q$が同等の主張であることを意味し、\textbf{$\ds p$と$\ds q$は同値である}と読む。\par
      この記号を模様して、主張$\ds P$を主張$\ds Q$で定義するを

      \begin{equation*}
        \begin{split}
          $\ds P \defineProposition Q$
        \end{split}
      \end{equation*}

      と書く。

      \paragraph{発展的な演算子}
        さらに発展的な演算子として次のものもある。

        \begin{equation*}
          \begin{split}
            $\ds p \veebar q$     & $\ds \defineProposition (p \land \lnot q) \lor (\lnot p \lor q)$\\
            $\ds p \uparrow q$    & $\ds \defineProposition \lnot p \lor \lnot q$\\
          \end{split}
        \end{equation*}
        
        これらは、上から順に\textbf{排他的論理和(xor)}、\textbf{否定論理積(nand)}と言い、主に電気回路、プログラミングで使われることが多い。
      
      \paragraph{恒真命題・恒偽命題}
        常に真・偽のどちらかを表す記号を以下に与える。

        \begin{equation*}
          \begin{split}
            $\ds \top$  & $\ds \defineProposition 常に真$\\
            $\ds \bot$  & $\ds \defineProposition 常に偽$\\
          \end{split}
        \end{equation*}

        これらを上から順に\textbf{恒真命題(トートロジー)}、\textbf{恒偽命題(コントラディクション)}と呼ぶ。

      \newpage

      \paragraph{よく用いられる恒真命題}
        以下に示す命題は、すべて常に成立する命題で、数学的議論をする上で大変便利でそれゆえ用いられる頻度が高いのでここで抑えておこう。

        \begin{center}
          \begin{table}[h]
            \centering
            \begin{tabular}{lll}

              & \begin{cases}
                $\ds p \land q \Longleftrightarrow q \land p\\
                $\ds p \lor q \Longleftrightarrow q \lor p\\
              \end{cases} & 交換則\\

              & \begin{cases}
                $\ds (p \land q) \land r \Longleftrightarrow p \land (q \land r)\\
                $\ds (p \lor q) \lor r \Longleftrightarrow p \lor (q \lor r)\\
              \end{cases} & 結合測\\

              & \begin{cases}
                $\ds p \land (q \lor r) \Longleftrightarrow (p \land q) \lor (p \land r)\\
                $\ds p \lor (q \land r) \Longleftrightarrow (p \lor q) \land (p \lor r)\\
              \end{cases} & 分配則\\

              & $\ds \lnot \lnot p \Longleftrightarrow p$ & 二重否定則\\

              & \begin{cases}
                $\ds \lnot (p \land q) \Longleftrightarrow \lnot p \lor \lnot q\\
                $\ds \lnot (p \lor q) \Longleftrightarrow \lnot p \land \lnot q\\
              \end{cases} & ド・モルガン(de Morgan)則\\

              & $\ds "p \Longrightarrow q" \Longleftrightarrow "\lnot p \Longleftarrow \lnot q"$ & 対偶則

            \end{tabular}
          \end{table}
        \end{center}
        
      \paragraph{ならば文の基礎定理}
        通称ならば文$\ds p \Longrightarrow q$は、真理表に与えられていることにならえば、少々感覚とずれるが次のように$\ds \land, \lor, \lnot$を用いて表現できる。

        \begin{center}
          \begin{table}[h]
            \centering
            \begin{tabular}{lll}

              $\ds "p \Longrightarrow q"$ & $\ds \Longleftrightarrow "\lnot p \lor (p \land q)"$                    & \\
                                          & $\ds \Longleftrightarrow "(\lnot p \lor p) \land (\lnot p \lor q)"$     & $\ds (\because 分配則)$\\
                                          & $\ds \Longleftrightarrow "\top \land (\lnot p \lor q)"$                 & $\ds (\because \lnot p \lor p は常に真)$\\
                                          & $\ds \Longleftrightarrow "\lnot p \lor q"$                              & 

            \end{tabular}
          \end{table}
        \end{center}
        
        \footnote{常に真である恒真命題$\ds "\lnot p \lor p" \Leftrightarrow \top$との論理積(and)は、これを省いても同じ主張であると言える}
        \footnote{\land で結ばれるものを\textbf{論理積}といい、\lor で結ばれるものを\textbf{論理和}という。}
        
        よってまとめれば

        \begin{equation*}
          \begin{split}
            $\ds "p \Longrightarrow q" \Longleftrightarrow "\lnot p \lor q"$
          \end{split}
        \end{equation*}

        である。
        
    \subsubsection{述語論理}

\end{document}
